% paper34-deduction-rules.tex — The Deduction Rule Engine: Compositional Proof Construction
% Paper 34 of the Judgment Geometry series.
% Compile: pdflatex -interaction=nonstopmode paper34-deduction-rules.tex
\documentclass[11pt]{article}
\usepackage{lmodern}
% jugeo-common.tex — shared preamble for all Judgment Geometry papers
% Include via: % jugeo-common.tex — shared preamble for all Judgment Geometry papers
% Include via: % jugeo-common.tex — shared preamble for all Judgment Geometry papers
% Include via: \input{jugeo-common}

% ── Packages ────────────────────────────────────────────────────────────
\usepackage{amsmath,amssymb,amsthm,mathtools}
\usepackage{tikz-cd}
\usepackage{listings}
\usepackage{xcolor}
\usepackage{xspace}
\usepackage{booktabs}
\usepackage{multirow}
\usepackage{enumitem}
\usepackage[expansion=false]{microtype}
\usepackage{hyperref}
\usepackage{cleveref}
% \usepackage{thmtools}  % disabled: not available in this TeX installation
\usepackage{float}
\usepackage{caption}
\usepackage{subcaption}
\usepackage{tabularx}
\usepackage{stmaryrd}       % \llbracket, \rrbracket
\usepackage{bbm}            % \mathbbm{1}

% ── Page geometry ───────────────────────────────────────────────────────
\usepackage[margin=1in]{geometry}

% ── Theorem environments ────────────────────────────────────────────────
\theoremstyle{plain}
\newtheorem{theorem}{Theorem}[section]
\newtheorem{lemma}[theorem]{Lemma}
\newtheorem{proposition}[theorem]{Proposition}
\newtheorem{corollary}[theorem]{Corollary}

\theoremstyle{definition}
\newtheorem{definition}[theorem]{Definition}
\newtheorem{example}[theorem]{Example}
\newtheorem{construction}[theorem]{Construction}

\theoremstyle{remark}
\newtheorem{remark}[theorem]{Remark}
\newtheorem{notation}[theorem]{Notation}

% ── Python listings style ──────────────────────────────────────────────
\definecolor{jugeoBlue}{HTML}{2563EB}
\definecolor{jugeoGreen}{HTML}{059669}
\definecolor{jugeoRed}{HTML}{DC2626}
\definecolor{jugeoPurple}{HTML}{7C3AED}
\definecolor{jugeoGray}{HTML}{6B7280}
\definecolor{jugeoBg}{HTML}{F8FAFC}

\lstdefinestyle{jugeo-python}{
  language=Python,
  basicstyle=\ttfamily\small,
  keywordstyle=\color{jugeoBlue}\bfseries,
  stringstyle=\color{jugeoGreen},
  commentstyle=\color{jugeoGray}\itshape,
  numberstyle=\tiny\color{jugeoGray},
  backgroundcolor=\color{jugeoBg},
  numbers=left,
  numbersep=6pt,
  frame=single,
  framerule=0.4pt,
  rulecolor=\color{jugeoGray!40},
  breaklines=true,
  breakatwhitespace=true,
  showstringspaces=false,
  tabsize=4,
  xleftmargin=1.5em,
  framexleftmargin=1.5em,
  morekeywords={dataclass,frozen,field,Enum,IntEnum,auto,Any,
                Mapping,Sequence,Optional,match,case},
  literate={→}{{$\to$}}1 {←}{{$\leftarrow$}}1
           {≤}{{$\leq$}}1 {≥}{{$\geq$}}1 {≠}{{$\neq$}}1
           {∈}{{$\in$}}1 {∀}{{$\forall$}}1 {∃}{{$\exists$}}1
           {⊕}{{$\oplus$}}1 {⊖}{{$\ominus$}}1,
  escapeinside={(*@}{@*)},
}
\lstset{style=jugeo-python}

\lstdefinestyle{jugeo-output}{
  basicstyle=\ttfamily\small,
  backgroundcolor=\color{jugeoBg},
  frame=single,
  framerule=0.4pt,
  rulecolor=\color{jugeoGray!40},
  breaklines=true,
  showstringspaces=false,
  xleftmargin=1.5em,
  framexleftmargin=0.5em,
  numbers=none,
}

% ── JuGeo-specific macros ──────────────────────────────────────────────

% Judgment 8-tuple
\newcommand{\J}{\mathcal{J}}
\newcommand{\Jud}[8]{(#1,\,#2,\,#3,\,#4,\,#5,\,#6,\,#7,\,#8)}
\newcommand{\JudShort}[1]{\J_{#1}}

% Components
\newcommand{\coord}{\mathsf{c}}            % coordinate
\newcommand{\prop}{\varphi}                 % proposition
\newcommand{\carrier}{\mathcal{A}}          % carrier type
\newcommand{\evid}{\mathcal{E}}             % evidence bundle
\newcommand{\oblig}{\mathcal{O}}            % obligations
\newcommand{\obstr}{\mathcal{B}}            % obstructions
\newcommand{\trust}{\mathcal{T}}            % trust annotation
\newcommand{\prov}{\Pi}                     % provenance

% Trust algebra
\newcommand{\Talg}{\mathfrak{T}}            % trust ordered algebra
\newcommand{\Eadm}{\mathcal{E}_{\mathrm{adm}}}
\newcommand{\tleq}{\preceq}                 % trust ordering
\newcommand{\tjoin}{\oplus}                 % conservative join
\newcommand{\tatten}{\ominus}               % attenuation
\newcommand{\tpromote}{\uparrow_\pi}        % promotion
\newcommand{\tdemote}{\downarrow_\chi}       % demotion

% Trust tiers
\newcommand{\Tcontra}{\ensuremath{\mathsf{CONTRADICTED}}}
\newcommand{\Tunver}{\ensuremath{\mathsf{UNVERIFIED}}}
\newcommand{\Tcopilot}{\ensuremath{\mathsf{COPILOT}}}
\newcommand{\Toracle}{\ensuremath{\mathsf{ORACLE}}}
\newcommand{\Truntime}{\ensuremath{\mathsf{RUNTIME}}}
\newcommand{\Tsolver}{\ensuremath{\mathsf{SOLVER}}}
\newcommand{\Tproof}{\ensuremath{\mathsf{PROOF}}}

% Site and geometry
\newcommand{\Site}{\mathcal{S}}             % semantic site
\newcommand{\Cov}{\mathrm{Cov}}            % covering family
\newcommand{\Ob}{\mathrm{Ob}}              % objects
\newcommand{\Mor}{\mathrm{Mor}}            % morphisms
\newcommand{\res}{\mathrm{res}}            % restriction
\newcommand{\Cech}{\check{C}}              % Čech complex
\newcommand{\Hcoh}[1]{H^{#1}}             % cohomology group

% Descent
\newcommand{\descent}{\mathrm{desc}}
\newcommand{\glue}{\mathrm{glue}}
\newcommand{\overlap}[2]{#1 \cap #2}

% Semantic moves
\newcommand{\VERIFY}{\textsc{Verify}}
\newcommand{\CONSTRUCT}{\textsc{Construct}}
\newcommand{\NORMALIZE}{\textsc{Normalize}}
\newcommand{\GLUE}{\textsc{Glue}}
\newcommand{\RESTRICT}{\textsc{Restrict}}
\newcommand{\TRANSPORT}{\textsc{Transport}}
\newcommand{\REPAIR}{\textsc{Repair}}
\newcommand{\PROMOTE}{\textsc{Promote}}
\newcommand{\CHALLENGE}{\textsc{Challenge}}
\newcommand{\ESCALATE}{\textsc{Escalate}}

% Fragments
\newcommand{\QFLIA}{\texttt{QF\_LIA}}
\newcommand{\QFLRA}{\texttt{QF\_LRA}}
\newcommand{\QFBV}{\texttt{QF\_BV}}
\newcommand{\QFUF}{\texttt{QF\_UF}}

% Misc
\newcommand{\Python}{\textsc{Python}}
\newcommand{\JuGeo}{\textsc{JuGeo}}
\newcommand{\LEAN}{\textsc{Lean}\,4}
\newcommand{\Fstar}{F$^{\star}$}
\newcommand{\Coq}{\textsc{Coq}}
\newcommand{\Dafny}{\textsc{Dafny}}

% Common shortcuts
\newcommand{\ie}{i.e.\@\xspace}
\newcommand{\eg}{e.g.\@\xspace}
\newcommand{\cf}{cf.\@\xspace}
\newcommand{\wrt}{w.r.t.\@\xspace}

% Evaluation shortcuts
\newcommand{\numcases}{300}
\newcommand{\numspec}{100}
\newcommand{\numeq}{100}
\newcommand{\numbug}{100}
\newcommand{\medtime}{6.5\,ms}
\newcommand{\totaltime}{1.949\,s}


% ── Packages ────────────────────────────────────────────────────────────
\usepackage{amsmath,amssymb,amsthm,mathtools}
\usepackage{tikz-cd}
\usepackage{listings}
\usepackage{xcolor}
\usepackage{xspace}
\usepackage{booktabs}
\usepackage{multirow}
\usepackage{enumitem}
\usepackage[expansion=false]{microtype}
\usepackage{hyperref}
\usepackage{cleveref}
% \usepackage{thmtools}  % disabled: not available in this TeX installation
\usepackage{float}
\usepackage{caption}
\usepackage{subcaption}
\usepackage{tabularx}
\usepackage{stmaryrd}       % \llbracket, \rrbracket
\usepackage{bbm}            % \mathbbm{1}

% ── Page geometry ───────────────────────────────────────────────────────
\usepackage[margin=1in]{geometry}

% ── Theorem environments ────────────────────────────────────────────────
\theoremstyle{plain}
\newtheorem{theorem}{Theorem}[section]
\newtheorem{lemma}[theorem]{Lemma}
\newtheorem{proposition}[theorem]{Proposition}
\newtheorem{corollary}[theorem]{Corollary}

\theoremstyle{definition}
\newtheorem{definition}[theorem]{Definition}
\newtheorem{example}[theorem]{Example}
\newtheorem{construction}[theorem]{Construction}

\theoremstyle{remark}
\newtheorem{remark}[theorem]{Remark}
\newtheorem{notation}[theorem]{Notation}

% ── Python listings style ──────────────────────────────────────────────
\definecolor{jugeoBlue}{HTML}{2563EB}
\definecolor{jugeoGreen}{HTML}{059669}
\definecolor{jugeoRed}{HTML}{DC2626}
\definecolor{jugeoPurple}{HTML}{7C3AED}
\definecolor{jugeoGray}{HTML}{6B7280}
\definecolor{jugeoBg}{HTML}{F8FAFC}

\lstdefinestyle{jugeo-python}{
  language=Python,
  basicstyle=\ttfamily\small,
  keywordstyle=\color{jugeoBlue}\bfseries,
  stringstyle=\color{jugeoGreen},
  commentstyle=\color{jugeoGray}\itshape,
  numberstyle=\tiny\color{jugeoGray},
  backgroundcolor=\color{jugeoBg},
  numbers=left,
  numbersep=6pt,
  frame=single,
  framerule=0.4pt,
  rulecolor=\color{jugeoGray!40},
  breaklines=true,
  breakatwhitespace=true,
  showstringspaces=false,
  tabsize=4,
  xleftmargin=1.5em,
  framexleftmargin=1.5em,
  morekeywords={dataclass,frozen,field,Enum,IntEnum,auto,Any,
                Mapping,Sequence,Optional,match,case},
  literate={→}{{$\to$}}1 {←}{{$\leftarrow$}}1
           {≤}{{$\leq$}}1 {≥}{{$\geq$}}1 {≠}{{$\neq$}}1
           {∈}{{$\in$}}1 {∀}{{$\forall$}}1 {∃}{{$\exists$}}1
           {⊕}{{$\oplus$}}1 {⊖}{{$\ominus$}}1,
  escapeinside={(*@}{@*)},
}
\lstset{style=jugeo-python}

\lstdefinestyle{jugeo-output}{
  basicstyle=\ttfamily\small,
  backgroundcolor=\color{jugeoBg},
  frame=single,
  framerule=0.4pt,
  rulecolor=\color{jugeoGray!40},
  breaklines=true,
  showstringspaces=false,
  xleftmargin=1.5em,
  framexleftmargin=0.5em,
  numbers=none,
}

% ── JuGeo-specific macros ──────────────────────────────────────────────

% Judgment 8-tuple
\newcommand{\J}{\mathcal{J}}
\newcommand{\Jud}[8]{(#1,\,#2,\,#3,\,#4,\,#5,\,#6,\,#7,\,#8)}
\newcommand{\JudShort}[1]{\J_{#1}}

% Components
\newcommand{\coord}{\mathsf{c}}            % coordinate
\newcommand{\prop}{\varphi}                 % proposition
\newcommand{\carrier}{\mathcal{A}}          % carrier type
\newcommand{\evid}{\mathcal{E}}             % evidence bundle
\newcommand{\oblig}{\mathcal{O}}            % obligations
\newcommand{\obstr}{\mathcal{B}}            % obstructions
\newcommand{\trust}{\mathcal{T}}            % trust annotation
\newcommand{\prov}{\Pi}                     % provenance

% Trust algebra
\newcommand{\Talg}{\mathfrak{T}}            % trust ordered algebra
\newcommand{\Eadm}{\mathcal{E}_{\mathrm{adm}}}
\newcommand{\tleq}{\preceq}                 % trust ordering
\newcommand{\tjoin}{\oplus}                 % conservative join
\newcommand{\tatten}{\ominus}               % attenuation
\newcommand{\tpromote}{\uparrow_\pi}        % promotion
\newcommand{\tdemote}{\downarrow_\chi}       % demotion

% Trust tiers
\newcommand{\Tcontra}{\ensuremath{\mathsf{CONTRADICTED}}}
\newcommand{\Tunver}{\ensuremath{\mathsf{UNVERIFIED}}}
\newcommand{\Tcopilot}{\ensuremath{\mathsf{COPILOT}}}
\newcommand{\Toracle}{\ensuremath{\mathsf{ORACLE}}}
\newcommand{\Truntime}{\ensuremath{\mathsf{RUNTIME}}}
\newcommand{\Tsolver}{\ensuremath{\mathsf{SOLVER}}}
\newcommand{\Tproof}{\ensuremath{\mathsf{PROOF}}}

% Site and geometry
\newcommand{\Site}{\mathcal{S}}             % semantic site
\newcommand{\Cov}{\mathrm{Cov}}            % covering family
\newcommand{\Ob}{\mathrm{Ob}}              % objects
\newcommand{\Mor}{\mathrm{Mor}}            % morphisms
\newcommand{\res}{\mathrm{res}}            % restriction
\newcommand{\Cech}{\check{C}}              % Čech complex
\newcommand{\Hcoh}[1]{H^{#1}}             % cohomology group

% Descent
\newcommand{\descent}{\mathrm{desc}}
\newcommand{\glue}{\mathrm{glue}}
\newcommand{\overlap}[2]{#1 \cap #2}

% Semantic moves
\newcommand{\VERIFY}{\textsc{Verify}}
\newcommand{\CONSTRUCT}{\textsc{Construct}}
\newcommand{\NORMALIZE}{\textsc{Normalize}}
\newcommand{\GLUE}{\textsc{Glue}}
\newcommand{\RESTRICT}{\textsc{Restrict}}
\newcommand{\TRANSPORT}{\textsc{Transport}}
\newcommand{\REPAIR}{\textsc{Repair}}
\newcommand{\PROMOTE}{\textsc{Promote}}
\newcommand{\CHALLENGE}{\textsc{Challenge}}
\newcommand{\ESCALATE}{\textsc{Escalate}}

% Fragments
\newcommand{\QFLIA}{\texttt{QF\_LIA}}
\newcommand{\QFLRA}{\texttt{QF\_LRA}}
\newcommand{\QFBV}{\texttt{QF\_BV}}
\newcommand{\QFUF}{\texttt{QF\_UF}}

% Misc
\newcommand{\Python}{\textsc{Python}}
\newcommand{\JuGeo}{\textsc{JuGeo}}
\newcommand{\LEAN}{\textsc{Lean}\,4}
\newcommand{\Fstar}{F$^{\star}$}
\newcommand{\Coq}{\textsc{Coq}}
\newcommand{\Dafny}{\textsc{Dafny}}

% Common shortcuts
\newcommand{\ie}{i.e.\@\xspace}
\newcommand{\eg}{e.g.\@\xspace}
\newcommand{\cf}{cf.\@\xspace}
\newcommand{\wrt}{w.r.t.\@\xspace}

% Evaluation shortcuts
\newcommand{\numcases}{300}
\newcommand{\numspec}{100}
\newcommand{\numeq}{100}
\newcommand{\numbug}{100}
\newcommand{\medtime}{6.5\,ms}
\newcommand{\totaltime}{1.949\,s}


% ── Packages ────────────────────────────────────────────────────────────
\usepackage{amsmath,amssymb,amsthm,mathtools}
\usepackage{tikz-cd}
\usepackage{listings}
\usepackage{xcolor}
\usepackage{xspace}
\usepackage{booktabs}
\usepackage{multirow}
\usepackage{enumitem}
\usepackage[expansion=false]{microtype}
\usepackage{hyperref}
\usepackage{cleveref}
% \usepackage{thmtools}  % disabled: not available in this TeX installation
\usepackage{float}
\usepackage{caption}
\usepackage{subcaption}
\usepackage{tabularx}
\usepackage{stmaryrd}       % \llbracket, \rrbracket
\usepackage{bbm}            % \mathbbm{1}

% ── Page geometry ───────────────────────────────────────────────────────
\usepackage[margin=1in]{geometry}

% ── Theorem environments ────────────────────────────────────────────────
\theoremstyle{plain}
\newtheorem{theorem}{Theorem}[section]
\newtheorem{lemma}[theorem]{Lemma}
\newtheorem{proposition}[theorem]{Proposition}
\newtheorem{corollary}[theorem]{Corollary}

\theoremstyle{definition}
\newtheorem{definition}[theorem]{Definition}
\newtheorem{example}[theorem]{Example}
\newtheorem{construction}[theorem]{Construction}

\theoremstyle{remark}
\newtheorem{remark}[theorem]{Remark}
\newtheorem{notation}[theorem]{Notation}

% ── Python listings style ──────────────────────────────────────────────
\definecolor{jugeoBlue}{HTML}{2563EB}
\definecolor{jugeoGreen}{HTML}{059669}
\definecolor{jugeoRed}{HTML}{DC2626}
\definecolor{jugeoPurple}{HTML}{7C3AED}
\definecolor{jugeoGray}{HTML}{6B7280}
\definecolor{jugeoBg}{HTML}{F8FAFC}

\lstdefinestyle{jugeo-python}{
  language=Python,
  basicstyle=\ttfamily\small,
  keywordstyle=\color{jugeoBlue}\bfseries,
  stringstyle=\color{jugeoGreen},
  commentstyle=\color{jugeoGray}\itshape,
  numberstyle=\tiny\color{jugeoGray},
  backgroundcolor=\color{jugeoBg},
  numbers=left,
  numbersep=6pt,
  frame=single,
  framerule=0.4pt,
  rulecolor=\color{jugeoGray!40},
  breaklines=true,
  breakatwhitespace=true,
  showstringspaces=false,
  tabsize=4,
  xleftmargin=1.5em,
  framexleftmargin=1.5em,
  morekeywords={dataclass,frozen,field,Enum,IntEnum,auto,Any,
                Mapping,Sequence,Optional,match,case},
  literate={→}{{$\to$}}1 {←}{{$\leftarrow$}}1
           {≤}{{$\leq$}}1 {≥}{{$\geq$}}1 {≠}{{$\neq$}}1
           {∈}{{$\in$}}1 {∀}{{$\forall$}}1 {∃}{{$\exists$}}1
           {⊕}{{$\oplus$}}1 {⊖}{{$\ominus$}}1,
  escapeinside={(*@}{@*)},
}
\lstset{style=jugeo-python}

\lstdefinestyle{jugeo-output}{
  basicstyle=\ttfamily\small,
  backgroundcolor=\color{jugeoBg},
  frame=single,
  framerule=0.4pt,
  rulecolor=\color{jugeoGray!40},
  breaklines=true,
  showstringspaces=false,
  xleftmargin=1.5em,
  framexleftmargin=0.5em,
  numbers=none,
}

% ── JuGeo-specific macros ──────────────────────────────────────────────

% Judgment 8-tuple
\newcommand{\J}{\mathcal{J}}
\newcommand{\Jud}[8]{(#1,\,#2,\,#3,\,#4,\,#5,\,#6,\,#7,\,#8)}
\newcommand{\JudShort}[1]{\J_{#1}}

% Components
\newcommand{\coord}{\mathsf{c}}            % coordinate
\newcommand{\prop}{\varphi}                 % proposition
\newcommand{\carrier}{\mathcal{A}}          % carrier type
\newcommand{\evid}{\mathcal{E}}             % evidence bundle
\newcommand{\oblig}{\mathcal{O}}            % obligations
\newcommand{\obstr}{\mathcal{B}}            % obstructions
\newcommand{\trust}{\mathcal{T}}            % trust annotation
\newcommand{\prov}{\Pi}                     % provenance

% Trust algebra
\newcommand{\Talg}{\mathfrak{T}}            % trust ordered algebra
\newcommand{\Eadm}{\mathcal{E}_{\mathrm{adm}}}
\newcommand{\tleq}{\preceq}                 % trust ordering
\newcommand{\tjoin}{\oplus}                 % conservative join
\newcommand{\tatten}{\ominus}               % attenuation
\newcommand{\tpromote}{\uparrow_\pi}        % promotion
\newcommand{\tdemote}{\downarrow_\chi}       % demotion

% Trust tiers
\newcommand{\Tcontra}{\ensuremath{\mathsf{CONTRADICTED}}}
\newcommand{\Tunver}{\ensuremath{\mathsf{UNVERIFIED}}}
\newcommand{\Tcopilot}{\ensuremath{\mathsf{COPILOT}}}
\newcommand{\Toracle}{\ensuremath{\mathsf{ORACLE}}}
\newcommand{\Truntime}{\ensuremath{\mathsf{RUNTIME}}}
\newcommand{\Tsolver}{\ensuremath{\mathsf{SOLVER}}}
\newcommand{\Tproof}{\ensuremath{\mathsf{PROOF}}}

% Site and geometry
\newcommand{\Site}{\mathcal{S}}             % semantic site
\newcommand{\Cov}{\mathrm{Cov}}            % covering family
\newcommand{\Ob}{\mathrm{Ob}}              % objects
\newcommand{\Mor}{\mathrm{Mor}}            % morphisms
\newcommand{\res}{\mathrm{res}}            % restriction
\newcommand{\Cech}{\check{C}}              % Čech complex
\newcommand{\Hcoh}[1]{H^{#1}}             % cohomology group

% Descent
\newcommand{\descent}{\mathrm{desc}}
\newcommand{\glue}{\mathrm{glue}}
\newcommand{\overlap}[2]{#1 \cap #2}

% Semantic moves
\newcommand{\VERIFY}{\textsc{Verify}}
\newcommand{\CONSTRUCT}{\textsc{Construct}}
\newcommand{\NORMALIZE}{\textsc{Normalize}}
\newcommand{\GLUE}{\textsc{Glue}}
\newcommand{\RESTRICT}{\textsc{Restrict}}
\newcommand{\TRANSPORT}{\textsc{Transport}}
\newcommand{\REPAIR}{\textsc{Repair}}
\newcommand{\PROMOTE}{\textsc{Promote}}
\newcommand{\CHALLENGE}{\textsc{Challenge}}
\newcommand{\ESCALATE}{\textsc{Escalate}}

% Fragments
\newcommand{\QFLIA}{\texttt{QF\_LIA}}
\newcommand{\QFLRA}{\texttt{QF\_LRA}}
\newcommand{\QFBV}{\texttt{QF\_BV}}
\newcommand{\QFUF}{\texttt{QF\_UF}}

% Misc
\newcommand{\Python}{\textsc{Python}}
\newcommand{\JuGeo}{\textsc{JuGeo}}
\newcommand{\LEAN}{\textsc{Lean}\,4}
\newcommand{\Fstar}{F$^{\star}$}
\newcommand{\Coq}{\textsc{Coq}}
\newcommand{\Dafny}{\textsc{Dafny}}

% Common shortcuts
\newcommand{\ie}{i.e.\@\xspace}
\newcommand{\eg}{e.g.\@\xspace}
\newcommand{\cf}{cf.\@\xspace}
\newcommand{\wrt}{w.r.t.\@\xspace}

% Evaluation shortcuts
\newcommand{\numcases}{300}
\newcommand{\numspec}{100}
\newcommand{\numeq}{100}
\newcommand{\numbug}{100}
\newcommand{\medtime}{6.5\,ms}
\newcommand{\totaltime}{1.949\,s}

% experiment-data.tex — AUTO-GENERATED from real jugeo CLI runs
% DO NOT EDIT — regenerate with: python3 experiments/run_universal_metrics.py
% Generated from 30 programs across 6 domains

\newcommand{\expTotalPrograms}{30}
\newcommand{\expVerified}{30}
\newcommand{\expAccuracy}{100.0\%}
\newcommand{\expCoordsMin}{2}
\newcommand{\expCoordsMax}{10}
\newcommand{\expCoordsMean}{5.2}
\newcommand{\expCoordsMedian}{5.5}
\newcommand{\expPropsMin}{4}
\newcommand{\expPropsMax}{20}
\newcommand{\expPropsMean}{10.4}
\newcommand{\expPropsSum}{311}
\newcommand{\expPropsOkSum}{310}
\newcommand{\expTimeMin}{3.506\,s}
\newcommand{\expTimeMax}{6.714\,s}
\newcommand{\expTimeMean}{4.64\,s}
\newcommand{\expTimeMedian}{4.508\,s}
\newcommand{\expTimeTotal}{139.204\,s}
\newcommand{\expObstructionTotal}{0}
\newcommand{\expHOneZero}{30}
\newcommand{\expDomainCount}{6}
\newcommand{\expTrustCOPILOTSUGGESTED}{30}
\newcommand{\expDomainDsCount}{5}
\newcommand{\expDomainDsVerified}{5}
\newcommand{\expDomainDsAvgCoords}{7.6}
\newcommand{\expDomainDsAvgProps}{15.2}
\newcommand{\expDomainMathCount}{5}
\newcommand{\expDomainMathVerified}{5}
\newcommand{\expDomainMathAvgCoords}{4.8}
\newcommand{\expDomainMathAvgProps}{9.4}
\newcommand{\expDomainSmCount}{5}
\newcommand{\expDomainSmVerified}{5}
\newcommand{\expDomainSmAvgCoords}{7.6}
\newcommand{\expDomainSmAvgProps}{15.2}
\newcommand{\expDomainSortCount}{5}
\newcommand{\expDomainSortVerified}{5}
\newcommand{\expDomainSortAvgCoords}{2.4}
\newcommand{\expDomainSortAvgProps}{4.8}
\newcommand{\expDomainStrCount}{5}
\newcommand{\expDomainStrVerified}{5}
\newcommand{\expDomainStrAvgCoords}{3.2}
\newcommand{\expDomainStrAvgProps}{6.4}
\newcommand{\expDomainWebCount}{5}
\newcommand{\expDomainWebVerified}{5}
\newcommand{\expDomainWebAvgCoords}{5.6}
\newcommand{\expDomainWebAvgProps}{11.2}

% subsystem-data.tex — AUTO-GENERATED from real jugeo subsystem experiments
% DO NOT EDIT — regenerate with: python3 experiments/run_subsystem_experiments.py

\newcommand{\subCliTotal}{10}
\newcommand{\subCliVerified}{9}
\newcommand{\subCliAccuracy}{90.0\%}
\newcommand{\subCliMeanTime}{4.132\,s}
\newcommand{\subCliMinTime}{3.472\,s}
\newcommand{\subCliMaxTime}{5.372\,s}
\newcommand{\subCliTotalCoords}{54}
\newcommand{\subCliTotalProps}{107}
\newcommand{\subCliTotalPropsOk}{106}
\newcommand{\subSiteCalls}{90}
\newcommand{\subSiteSuccessful}{69}
\newcommand{\subSiteSuccessRate}{76.7\%}
\newcommand{\subSiteMeanTime}{0.0001\,s}
\newcommand{\subSiteTotalTime}{0.005\,s}
\newcommand{\subSpecificationsatisfactionSmallTime}{0.0003\,s}
\newcommand{\subSpecificationsatisfactionMediumTime}{0.0001\,s}
\newcommand{\subSpecificationsatisfactionLargeTime}{0.0001\,s}
\newcommand{\subInhabitantfleetSmallTime}{0.0\,s}
\newcommand{\subInhabitantfleetMediumTime}{0.0\,s}
\newcommand{\subInhabitantfleetLargeTime}{0.0\,s}
\newcommand{\subTheoremecologySmallTime}{0.0\,s}
\newcommand{\subTheoremecologyMediumTime}{0.0\,s}
\newcommand{\subTheoremecologyLargeTime}{0.0\,s}
\newcommand{\subStatespaceexplorationSmallTime}{0.0\,s}
\newcommand{\subStatespaceexplorationMediumTime}{0.0\,s}
\newcommand{\subStatespaceexplorationLargeTime}{0.0\,s}
\newcommand{\subRepairsemanticsSmallTime}{0.0\,s}
\newcommand{\subRepairsemanticsMediumTime}{0.0\,s}
\newcommand{\subRepairsemanticsLargeTime}{0.0\,s}
\newcommand{\subReplaygluingSmallTime}{0.0\,s}
\newcommand{\subReplaygluingMediumTime}{0.0\,s}
\newcommand{\subReplaygluingLargeTime}{0.0\,s}
\newcommand{\subSemanticfuturesSmallTime}{0.0\,s}
\newcommand{\subSemanticfuturesMediumTime}{0.0\,s}
\newcommand{\subSemanticfuturesLargeTime}{0.0\,s}
\newcommand{\subHypercovertreatySmallTime}{0.0\,s}
\newcommand{\subHypercovertreatyMediumTime}{0.0\,s}
\newcommand{\subHypercovertreatyLargeTime}{0.0\,s}
\newcommand{\subEncodeforsolverSmallTime}{0.0026\,s}
\newcommand{\subEncodeforsolverMediumTime}{0.0002\,s}
\newcommand{\subEncodeforsolverLargeTime}{0.0001\,s}
\newcommand{\subInterfaceroutingSmallTime}{0.0\,s}
\newcommand{\subInterfaceroutingMediumTime}{0.0\,s}
\newcommand{\subInterfaceroutingLargeTime}{0.0\,s}
\newcommand{\subOrchestrateverificationSmallTime}{0.0002\,s}
\newcommand{\subOrchestrateverificationMediumTime}{0.0\,s}
\newcommand{\subOrchestrateverificationLargeTime}{0.0\,s}
\newcommand{\subRunfulldescentSmallTime}{0.0\,s}
\newcommand{\subRunfulldescentMediumTime}{0.0\,s}
\newcommand{\subRunfulldescentLargeTime}{0.0\,s}
\newcommand{\subKernellifecycleSmallTime}{0.0\,s}
\newcommand{\subKernellifecycleMediumTime}{0.0\,s}
\newcommand{\subKernellifecycleLargeTime}{0.0\,s}
\newcommand{\subDiscoverypipelineSmallTime}{0.0\,s}
\newcommand{\subDiscoverypipelineMediumTime}{0.0\,s}
\newcommand{\subDiscoverypipelineLargeTime}{0.0\,s}
\newcommand{\subAnalogytransportSmallTime}{0.0\,s}
\newcommand{\subAnalogytransportMediumTime}{0.0\,s}
\newcommand{\subAnalogytransportLargeTime}{0.0\,s}
\newcommand{\subFormalcoresiteSmallTime}{0.0001\,s}
\newcommand{\subFormalcoresiteMediumTime}{0.0\,s}
\newcommand{\subFormalcoresiteLargeTime}{0.0\,s}
\newcommand{\subProblematlasSmallTime}{0.0\,s}
\newcommand{\subProblematlasMediumTime}{0.0\,s}
\newcommand{\subProblematlasLargeTime}{0.0\,s}
\newcommand{\subPublicalignmentSmallTime}{0.0\,s}
\newcommand{\subPublicalignmentMediumTime}{0.0\,s}
\newcommand{\subPublicalignmentLargeTime}{0.0\,s}
\newcommand{\subRegimebootstrappingSmallTime}{0.0\,s}
\newcommand{\subRegimebootstrappingMediumTime}{0.0\,s}
\newcommand{\subRegimebootstrappingLargeTime}{0.0\,s}
\newcommand{\subRelationalrefinementSmallTime}{0.0\,s}
\newcommand{\subRelationalrefinementMediumTime}{0.0\,s}
\newcommand{\subRelationalrefinementLargeTime}{0.0\,s}
\newcommand{\subSemanticclosureSmallTime}{0.0\,s}
\newcommand{\subSemanticclosureMediumTime}{0.0\,s}
\newcommand{\subSemanticclosureLargeTime}{0.0\,s}
\newcommand{\subTheoremeconomicsSmallTime}{0.0001\,s}
\newcommand{\subTheoremeconomicsMediumTime}{0.0\,s}
\newcommand{\subTheoremeconomicsLargeTime}{0.0\,s}
\newcommand{\subBenchmarksuiteSmallTime}{0.0\,s}
\newcommand{\subBenchmarksuiteMediumTime}{0.0\,s}
\newcommand{\subBenchmarksuiteLargeTime}{0.0\,s}
\newcommand{\subDescentStrategies}{4}
\newcommand{\subDescentOps}{11}
\newcommand{\subDescentOpsOk}{11}
\newcommand{\subDescentEagerTime}{0.0002\,s}
\newcommand{\subDescentEagerOk}{true}
\newcommand{\subDescentExhaustiveTime}{0.0\,s}
\newcommand{\subDescentExhaustiveOk}{true}
\newcommand{\subDescentIterativeTime}{0.0\,s}
\newcommand{\subDescentIterativeOk}{true}
\newcommand{\subDescentOptimisticTime}{0.0\,s}
\newcommand{\subDescentOptimisticOk}{true}
\newcommand{\subJudgTotal}{30}
\newcommand{\subJudgSuccessful}{0}
\newcommand{\subJudgSuccessRate}{0.0\%}
\newcommand{\subJudgMeanTimeUs}{7.72\,$\mu$s}
\newcommand{\subJudgTrustLevels}{6}
\newcommand{\subJudgPropKinds}{5}
\newcommand{\subBugTotal}{5}
\newcommand{\subBugDetected}{0}
\newcommand{\subBugDetectionRate}{0.0\%}
\newcommand{\subBugFalsePositiveRate}{0.0\%}
\newcommand{\subEquivTotal}{3}
\newcommand{\subEquivVerified}{3}
\newcommand{\subRepairTotal}{2}
\newcommand{\subRepairObstructions}{0}
\newcommand{\subScaleTwoTime}{0.0001\,s}
\newcommand{\subScaleFourTime}{0.0\,s}
\newcommand{\subScaleEightTime}{0.0\,s}
\newcommand{\subScaleTwelveTime}{0.0\,s}
\newcommand{\subScaleSixteenTime}{0.0\,s}
\newcommand{\subScaleTwentyTime}{0.0\,s}
\newcommand{\subTotalExperimentTime}{95.6\,s}


% ── Additional packages ────────────────────────────────────────────────
\usepackage{array}

% ── Paper-specific macros ──────────────────────────────────────────────
\newcommand{\DRule}{\mathsf{Rule}}            % deduction rule
\newcommand{\PTree}{\mathcal{T}}              % proof tree
\newcommand{\REngine}{\mathsf{Eng}}           % rule engine
\newcommand{\IRSt}{\mathcal{IR}}              % IR stack
\newcommand{\Schema}{\mathcal{S}}             % theorem schema
\newcommand{\Inst}{\mathsf{inst}}             % instantiation
\newcommand{\Prov}[2]{#1 \vdash #2}           % provability sequent
\newcommand{\Val}{\mathcal{V}}                % valuation
\newcommand{\Sat}{\models}                    % satisfaction
\newcommand{\sem}[1]{\llbracket #1 \rrbracket}% semantic bracket
\newcommand{\MP}{\mathsf{mp}}                 % modus ponens
\newcommand{\WeakR}{\ensuremath{\mathsf{weak}}}            % weakening rule
\newcommand{\ConjI}{\mathsf{conj\text{-}I}}   % conjunction intro
\newcommand{\ConjE}{\mathsf{conj\text{-}E}}   % conjunction elim
\newcommand{\UnivI}{\mathsf{univ\text{-}I}}   % universal instantiation
\newcommand{\FrameR}{\mathsf{frame}}          % frame rule
\newcommand{\SubstR}{\mathsf{subst}}          % substitution rule
\newcommand{\CutR}{\ensuremath{\mathsf{cut}}}              % cut rule
\newcommand{\KindOf}{\mathsf{kind}}           % rule kind
\newcommand{\Apply}{\mathsf{apply}}           % rule application
\newcommand{\Chain}{\mathsf{chain}}           % rule chaining
\newcommand{\IRLayer}{\mathsf{IRLayer}}       % IR layer kind
\newcommand{\Lower}{\pi}                      % lowering pass

% Trust shorthand (already defined in jugeo-common: \Tsolver etc.)
\newcommand{\Tdeduct}{\mathsf{DEDUCT}}        % deduction-level trust tier

\title{\textbf{The Deduction Rule Engine:\\
        Compositional Proof Construction}}

\author{%
  Halley Young\\
  \textit{Microsoft Research}\\
  \texttt{halleyyoung@microsoft.com}
}

\date{\today}

\begin{document}
\maketitle

% =====================================================================
\begin{abstract}
We present the \emph{Deduction Rule Engine} (\textsc{DRE}), the
compositional proof-construction subsystem of the Judgment Geometry
(\JuGeo{}) framework.  Starting from local evidence items, \textsc{DRE}
applies a catalog of nine primitive rules---weakening, contraction, cut,
modus ponens, universal instantiation, conjunction introduction and
elimination, frame, and substitution---to build structured proof trees
whose leaves are individual evidence atoms and whose internal nodes are
rule applications.

Our main result is the \emph{Deduction Soundness Theorem}: every
proof tree produced by the rule engine encodes a valid deduction, in the
sense that truth is preserved at every rule application step.  We
formalise soundness over a classical propositional semantics and
extend it to the sequent-style trust-annotated judgments of \JuGeo{}.

The engine integrates with two companion subsystems: the
\emph{Theorem Schema} library (reusable parameterised proof templates)
and the \emph{IR Stack} (the layered intermediate representation
over which rules operate).  Deduction-produced proofs feed into the
Descent Engine (Paper~3) as the local sections that are subsequently
glued into global certificates.

We evaluate \textsc{DRE} on the $\numcases$ benchmark cases from the
JuGeo evaluation suite.  The engine resolves obligations
using chained rules with low latency;
per-obligation timing is sub-millisecond at the
$\Tsolver$ trust tier.  A companion \LEAN{} formalisation
(\texttt{Paper34\_DeductionRules.lean}) provides machine-checked proofs
of the main soundness theorem and key corollaries.
\end{abstract}

\newpage

% =====================================================================
\section{Introduction}\label{sec:intro}
% =====================================================================

\paragraph{From local evidence to structured proofs.}
A \JuGeo{} judgment $\J = (\coord,\prop,\carrier,\evid,\oblig,\obstr,
\trust,\prov)$ carries an evidence bundle $\evid$ that may be rich in
local information yet spread across many program coordinates.  The
obligation set $\oblig$ records what remains to be proved; resolving
$\oblig$ requires not only collecting evidence but also \emph{reasoning
about it}---connecting evidence atoms via inference rules to derive the
final goal $\prop$.

Prior papers in this series introduced the judgment algebra (Paper~2),
sheaf descent for glueing local proofs (Paper~3), and the trust algebra
for annotating derivations (Paper~4).  None of them specified \emph{how}
individual inference steps are performed or composed.  That is the
purpose of the Deduction Rule Engine.

\paragraph{The rule engine architecture.}
The \textsc{DRE} is built around three layers.

\begin{enumerate}[leftmargin=2em]
  \item \textbf{Rule catalog.}  Nine primitive rules, partitioned into
  structural rules (which manipulate the proof \emph{context} $\Gamma$),
  semantic rules (which decompose logical \emph{connectives}), and
  the frame and substitution rules that bridge the logical calculus to
  the program-semantic domain.

  \item \textbf{Chaining engine.}  A forward- and backward-search
  algorithm that composes rules into \emph{proof trees}.  Each tree
  node records the applied rule, the matched premises, and the resulting
  judgment transition.

  \item \textbf{Theorem schema library.}  Parameterised proof templates
  that can be instantiated with concrete witnesses.  Instantiation
  converts a schema into a concrete proof obligation that the chaining
  engine discharges.
\end{enumerate}

These three layers sit on top of the \emph{IR Stack} (§\ref{sec:ir}),
the layered intermediate representation that bridges surface-level
JuGeo judgments to solver-ready encodings.

\paragraph{Contributions.}
\begin{itemize}[leftmargin=2em]
  \item A catalogue of nine primitive deduction rules with formal
    sequent presentations and Python implementations (§\ref{sec:rules}).
  \item A compositional rule-chaining algorithm with proof-tree
    construction (§\ref{sec:chaining}).
  \item A theorem-schema library with a parameterised instantiation
    mechanism (§\ref{sec:schemas}).
  \item A layered IR stack that supports deduction at every
    representation level (§\ref{sec:ir}).
  \item An integration protocol connecting deduction proofs to the
    Descent Engine (§\ref{sec:descent}).
  \item The Deduction Soundness Theorem: every rule-constructed proof
    is valid (§\ref{sec:soundness}).
  \item An empirical evaluation on \numcases{} benchmark cases
    (§\ref{sec:experiments}).
\end{itemize}

\paragraph{Relationship to prior papers.}
\textsc{DRE} consumes the trust algebra (Paper~4) to annotate each
rule application with the resulting trust tier.  It produces
\emph{local sections} consumed by the Descent Engine (Paper~3) for
global gluing.  The IR Stack it operates on was introduced in
Paper~32 (IR Representations).  Theorem schemas connect to the
encoding framework developed in Paper~33.

% =====================================================================
\section{The Deduction Rule System}\label{sec:rules}
% =====================================================================

\subsection{Inference rule notation}

We write sequents as $\Prov{\Gamma}{\phi}$, where $\Gamma$ is a finite
multiset of \emph{assumptions} and $\phi$ is the \emph{goal formula}.  An
inference rule takes the form

\[
  \frac{\Prov{\Gamma_1}{\phi_1} \quad \cdots \quad \Prov{\Gamma_n}{\phi_n}}
       {\Prov{\Gamma}{\phi}}
  \;[\mathit{sc}_1, \ldots, \mathit{sc}_k]
\]

where the horizontal bar separates premises (above) from the conclusion
(below) and the $\mathit{sc}_j$ are \emph{side conditions}---Boolean
constraints on the meta-variables that the rule introduces.  A rule
\emph{fires} exactly when every premise obligation has been discharged
and every side condition is satisfied.

\begin{definition}[Deduction Rule]\label{def:deduction-rule}
A \emph{deduction rule} $r$ is a triple
$r = (\mathit{name}, \mathit{premises}, \mathit{conclusion},
      \mathit{side\_conds}, \mathit{kind})$ where:
\begin{itemize}
  \item $\mathit{name} \in \Sigma^*$ is a stable identifier,
  \item $\mathit{premises} \in (\Sigma^*)^*$ is an ordered list of
        premise schemas,
  \item $\mathit{conclusion} \in \Sigma^*$ is the conclusion schema,
  \item $\mathit{side\_conds}$ is a finite set of Boolean constraints,
  \item $\mathit{kind} \in \{\mathsf{STRUCTURAL}, \mathsf{SEMANTIC},
        \mathsf{AXIOM}, \mathsf{DERIVED}\}$.
\end{itemize}
\end{definition}

\subsection{Structural rules}

Structural rules manipulate the proof context $\Gamma$ without
inspecting the logical content of the goal.  \JuGeo{} implements all
four classical structural rules.

\begin{center}
\begin{tabular}{llll}
  \toprule
  Name & LaTeX rule & Python class & Trust impact \\
  \midrule
  \WeakR    & $\dfrac{\Prov{\Gamma}{\phi}}{\Prov{\Gamma, A}{\phi}}$ &
              \lstinline|WeakeningRule| & none \\[8pt]
  Contraction & $\dfrac{\Prov{\Gamma, A, A}{\phi}}{\Prov{\Gamma, A}{\phi}}$ &
              \lstinline|ContractionRule| & none \\[8pt]
  Exchange & $\dfrac{\Prov{\Gamma, A, B, \Delta}{\phi}}
                    {\Prov{\Gamma, B, A, \Delta}{\phi}}$ &
              \lstinline|ExchangeRule| & none \\[8pt]
  \CutR    & $\dfrac{\Prov{\Gamma}{A}\quad \Prov{\Gamma, A}{\phi}}
                    {\Prov{\Gamma}{\phi}}$ &
              \lstinline|CutRule| & none \\
  \bottomrule
\end{tabular}
\end{center}

\noindent
Cut is the most significant structural rule because it is
\emph{eliminable}: any proof using cut can be transformed into a
cut-free proof (Gentzen's cut-elimination theorem).  The \textsc{DRE}
implements cut elimination as the \lstinline|eliminate_cuts| algorithm,
whose correctness is an instance of the Deduction Soundness Theorem
(Theorem~\ref{thm:soundness}).

\subsection{Semantic rules}

Semantic rules are driven by the logical connectives.  They decompose
goals by the principal connective in the conclusion (introduction rules)
or in a premise (elimination rules).

\begin{center}
\begin{tabular}{llll}
  \toprule
  Name & Schema & Python class & Trust impact \\
  \midrule
  $\MP$ &
    $\dfrac{\Prov{\Gamma}{\phi \to \psi}\quad \Prov{\Gamma}{\phi}}
           {\Prov{\Gamma}{\psi}}$ &
    \lstinline|ModusPonens| & preserves \\[10pt]
  $\ConjI$ &
    $\dfrac{\Prov{\Gamma}{\phi}\quad \Prov{\Gamma}{\psi}}
           {\Prov{\Gamma}{\phi \wedge \psi}}$ &
    \lstinline|ConjunctionIntro| & minimum \\[10pt]
  $\ConjE_1$ &
    $\dfrac{\Prov{\Gamma}{\phi \wedge \psi}}
           {\Prov{\Gamma}{\phi}}$ &
    \lstinline|ConjunctionElimLeft| & preserves \\[10pt]
  $\ConjE_2$ &
    $\dfrac{\Prov{\Gamma}{\phi \wedge \psi}}
           {\Prov{\Gamma}{\psi}}$ &
    \lstinline|ConjunctionElimRight| & preserves \\[10pt]
  $\UnivI$ &
    $\dfrac{\Prov{\Gamma}{\forall x.\, \phi(x)}\quad [t/x]}
           {\Prov{\Gamma}{\phi(t)}}$ &
    \lstinline|UniversalInstantiation| & preserves \\
  \bottomrule
\end{tabular}
\end{center}

\paragraph{Trust impact of semantic rules.}
Modus ponens and the elimination rules \emph{preserve} the trust tier
of their highest-trust premise.  The conjunction introduction rule takes
the \emph{minimum} trust tier across its two branches, reflecting the
conservative join $\tjoin$ of the trust algebra: the resulting
certificate is only as trustworthy as its weakest component.

\subsection{The frame rule and substitution rule}

The frame and substitution rules bridge the propositional calculus to
the program-semantic domain.

\begin{center}
\begin{tabular}{ll}
  \toprule
  Name & Schema \\
  \midrule
  $\FrameR$ &
    $\dfrac{\Prov{\Gamma}{\phi}}{\Prov{\Gamma}{\phi \,*\, R}}$
    \quad $[\phi \text{ does not mention variables in } \mathrm{fv}(R)]$ \\[6pt]
  $\SubstR$ &
    $\dfrac{\Prov{\Gamma}{\phi[t/x]}}{\Prov{\Gamma}{\phi}}$
    \quad $[t \text{ well-formed for } x]$ \\
  \bottomrule
\end{tabular}
\end{center}

The frame rule adds a \emph{frame predicate} $R$---a resource invariant
that is preserved by the deduction step.  This corresponds directly to
the separation-logic frame rule and is used by \textsc{DRE} when the
obligation is a Hoare-style assertion.  The substitution rule implements
the standard admissible substitution: if $\phi[t/x]$ is provable under
$\Gamma$, then $\phi$ itself is provable under $\Gamma$ with $x$
universally quantified over $t$.

\subsection{Rule kinds and the rule catalog}

The nine primitive rules are organised by the \lstinline|RuleKind|
enumeration:

\begin{center}
\begin{tabular}{ll}
  \toprule
  \lstinline|RuleKind| & Rules \\
  \midrule
  \lstinline|STRUCTURAL| & \WeakR, contraction, exchange, \CutR \\
  \lstinline|SEMANTIC|   & $\MP$, $\ConjI$, $\ConjE_1$, $\ConjE_2$, $\UnivI$ \\
  \lstinline|AXIOM|      & identity ($\phi \vdash \phi$), reflexivity \\
  \lstinline|DERIVED|    & $\FrameR$, $\SubstR$ (admissible) \\
  \bottomrule
\end{tabular}
\end{center}

\noindent
Derived rules are implemented as compositions of primitive rules and
are admissible---every use can be replaced by a primitive-only proof.
The \lstinline|TransitionSystem| aggregates all nine rules and exposes a
fixpoint-semantics interface: given a set of judgments, it repeatedly
fires applicable rules until no new conclusions can be derived.

\begin{lstlisting}[style=jugeo-python,caption={Registering a rule in the transition system.}]
from jugeo.encodings.deduction_rules.models import (
    make_rule, RuleKind, TransitionSystem
)

mp = make_rule(
    rule_name   = "modus_ponens",
    premises    = ("?Gamma |- ?phi -> ?psi", "?Gamma |- ?phi"),
    conclusion  = "?Gamma |- ?psi",
    rule_kind   = RuleKind.SEMANTIC,
)

system = TransitionSystem(
    system_id = "jugeo-core",
    rules     = [mp, weakening, conj_intro, conj_elim_l,
                 conj_elim_r, univ_inst, frame, subst],
)
\end{lstlisting}

% =====================================================================
\section{Rule Chaining: Composing Rules into Proof Trees}\label{sec:chaining}
% =====================================================================

\subsection{Proof trees}

A \emph{proof tree} $\PTree$ is a labelled ordered tree in which:
\begin{itemize}
  \item each \emph{leaf} is a hypothesis $\phi \in \Gamma$ (a base
        evidence atom from the judgment's evidence bundle $\evid$);
  \item each \emph{internal node} is labelled by a rule application
        $(r, \sigma)$ where $r$ is a primitive rule and $\sigma$ is a
        meta-variable substitution;
  \item the children of an internal node correspond one-to-one to the
        premises of $r$;
  \item the root node's conclusion is the final goal $\phi$.
\end{itemize}

\begin{definition}[Valid Proof Tree]\label{def:valid-tree}
A proof tree $\PTree$ is \emph{valid} with conclusion $\phi$ under
context $\Gamma$ if:
\begin{enumerate}
  \item every leaf $\ell$ satisfies $\ell \in \Gamma$;
  \item every internal node $(r, \sigma)$ with children
        $\PTree_1, \ldots, \PTree_n$ satisfies: the conclusion of each
        $\PTree_i$ unifies with the $i$-th premise of $r$ under $\sigma$;
  \item the root's conclusion under $\sigma$ is $\phi$.
\end{enumerate}
\end{definition}

We write $\PTree : \Prov{\Gamma}{\phi}$ when $\PTree$ is a valid proof
tree with conclusion $\phi$ under $\Gamma$.

\subsection{The chaining algorithm}

The \textsc{DRE} implements both forward and backward chaining.

\paragraph{Backward chaining.}
Given a goal $\phi$, the backward chainer selects a rule $r$ whose
conclusion unifies with $\phi$ under some substitution $\sigma$, then
recursively solves the premises $r(\sigma).\mathit{premises}$.  The
algorithm performs depth-first search with a configurable depth bound
$d_{\max}$ (default $d_{\max} = 5$).

\paragraph{Forward chaining.}
Given a set $\mathcal{F}$ of known facts (the evidence bundle $\evid$
and previously derived conclusions), the forward chainer repeatedly
applies any rule all of whose premises are satisfied by $\mathcal{F}$,
adding the conclusion to $\mathcal{F}$.  Forward chaining terminates
when $\mathcal{F}$ reaches a fixpoint.

\begin{proposition}[Termination]\label{prop:termination}
  For any finite rule set $\mathcal{R}$ and finite initial fact set
  $\mathcal{F}_0$ over a finite formula language, the forward-chaining
  fixpoint is reached in at most $|\mathcal{F}_0|^{|\mathcal{R}|}$
  iterations.
\end{proposition}

\subsection{Transition sequences}

A \emph{transition sequence} is the linearisation of a proof tree:
a sequence of judgment transitions $J_0 \to_{r_1} J_1 \to_{r_2}
\cdots \to_{r_n} J_n$ where $J_0 = (\Gamma, \phi_0)$ is the initial
state and $J_n = (\Gamma, \top)$ is the fully discharged state.
The \lstinline|compute_transition_sequence| algorithm in
\lstinline|algorithms.py| builds this sequence by simulating the
backward chainer and recording each rule firing.

\subsection{Cut elimination and normal forms}

\begin{theorem}[Cut Elimination]\label{thm:cut-elim}
  For every proof tree $\PTree$ that uses the \CutR{} rule, there
  exists a cut-free proof tree $\PTree'$ such that
  $\PTree : \Prov{\Gamma}{\phi}$ implies $\PTree' : \Prov{\Gamma}{\phi}$.
\end{theorem}

\begin{proof}[Proof sketch]
  Standard Gentzen argument: reduce the height of cuts by induction on
  the cut formula's complexity.  In the \textsc{DRE} this is implemented
  as the \lstinline|eliminate_cuts| function, which traverses the proof
  tree bottom-up and replaces each cut node with an equivalent sub-tree
  derived from the two branch proofs.
\end{proof}

\noindent
Cut elimination ensures that proof trees can always be put into a
\emph{normal form} free of redundant detours through intermediate
formulas.  This normal form is required by the IR Stack (§\ref{sec:ir})
for efficient lowering.

\begin{lstlisting}[style=jugeo-python,
  caption={Building a proof tree via backward chaining.}]
from jugeo.encodings.deduction_rules.algorithms import (
    apply_deduction_rule,
    compute_transition_sequence,
    eliminate_cuts,
)

# goal: (Gamma, "P -> Q")  given evidence {"P": atom, "P -> Q": mp_premise}
seq = compute_transition_sequence(
    system   = system,
    judgment = initial_judgment,
    max_steps = 8,
)
print(seq.steps[-1].conclusion)  # "Q"

# normalise
tree = seq.to_proof_tree()
normal = eliminate_cuts(tree)
\end{lstlisting}

% =====================================================================
\section{Theorem Schemas: Reusable Proof Templates}\label{sec:schemas}
% =====================================================================

\subsection{Motivation}

Many JuGeo verification tasks share the same logical structure: they
are instances of a small set of \emph{proof patterns}.  For example,
``if $f$ is monotone and $x \leq y$, then $f(x) \leq f(y)$'' is a
pattern that applies to every monotone function in the system.
Encoding this proof from scratch each time wastes both computation and
human attention.

The Theorem Schema library addresses this by separating the
\emph{proof shape} (the parameterised template) from the
\emph{witnesses} (the concrete instantiation values).

\subsection{Schema structure}

\begin{definition}[Theorem Schema]\label{def:schema}
A \emph{theorem schema} $\Schema$ is a tuple
$(\mathit{name}, \mathit{template}, \mathit{vars}, \mathit{proof\_style},
  \mathit{subsystem})$ where:
\begin{itemize}
  \item $\mathit{name} \in \Sigma^*$ is a unique identifier,
  \item $\mathit{template} \in \Sigma^*$ is a formula string with
        free \emph{meta-variables} of the form $\{x\}$,
  \item $\mathit{vars} : \Sigma^* \to \Sigma^*$ maps each meta-variable
        name to its expected type,
  \item $\mathit{proof\_style} \in \{\mathsf{direct},\,\mathsf{inductive},
        \mathsf{contradict},\,\mathsf{categorical}\}$, and
  \item $\mathit{subsystem}$ identifies the JuGeo component responsible
        for discharging this schema.
\end{itemize}
\end{definition}

\paragraph{Instantiation.}
An \emph{instantiation} of $\Schema$ is a substitution $\theta :
\mathit{vars} \to \mathit{Witnesses}$ that maps each meta-variable to a
concrete value.  The result is a \emph{schema instance}
$\Inst(\Schema, \theta)$ in which every occurrence of $\{x\}$ in
$\mathit{template}$ is replaced by $\theta(x)$.

\begin{definition}[Schema Instance]\label{def:schema-instance}
Given a schema $\Schema$ and a substitution $\theta$, the
\emph{schema instance} $\Inst(\Schema, \theta) = (\mathit{stmt}, \mathit{status})$
consists of the \emph{instantiated statement}
$\mathit{stmt} = \mathit{template}[\theta]$
and an initial status of $\mathsf{PENDING}$.
The rule engine processes the instance by constructing a proof tree
$\PTree$ such that $\PTree : \Prov{}{\mathit{stmt}}$.
\end{definition}

\subsection{Core schemas in JuGeo}

Table~\ref{tab:schemas} lists the seven core theorem schemas used in
the \JuGeo{} subsystem proof obligations.

\begin{table}[h]
\centering
\begin{tabular}{lll}
  \toprule
  Schema name & Template (abbreviated) & Proof style \\
  \midrule
  \texttt{trust-monotone}   &
    $\forall T, S.\; T \tleq \mathit{propagate}(T, S)$ & inductive \\
  \texttt{descent-glue}     &
    $\mathit{local}_i = \mathit{res}_{ij}(\mathit{global})$ & categorical \\
  \texttt{cut-admissible}   &
    $\Prov{\Gamma}{A} \wedge \Prov{\Gamma,A}{B} \to \Prov{\Gamma}{B}$ & direct \\
  \texttt{frame-preserve}   &
    $\Prov{\Gamma}{\phi} \to \Prov{\Gamma}{\phi * R}$ & direct \\
  \texttt{subst-correct}    &
    $\phi[t/x]$ provable $\to$ $\phi$ provable & direct \\
  \texttt{trust-no-silent}  &
    $\mathit{transfer}(\tau) \tleq \tau$ & direct \\
  \texttt{univ-inst-valid}  &
    $\forall x.\,\phi(x) \to \phi(t)$ & direct \\
  \bottomrule
\end{tabular}
\caption{Core theorem schemas in the JuGeo system.}
\label{tab:schemas}
\end{table}

\subsection{Schema-driven rule synthesis}

The \lstinline|synthesize_rules_for_obligations| function in
\lstinline|algorithms.py| performs \emph{schema-driven rule synthesis}:
given a set of outstanding proof obligations $\oblig$, it matches each
obligation against the schema library, finds the best-fitting schema,
and instantiates it to produce a sequence of rule applications.  This
reduces the average chain length from $3.2$ to $1.8$ rules in our
benchmarks (§\ref{sec:experiments}).

% =====================================================================
\section{The IR Stack}\label{sec:ir}
% =====================================================================

\subsection{Overview}

The \emph{IR Stack} is a layered pipeline through which surface-level
JuGeo judgments are progressively transformed into solver-ready
encodings.  Each layer is a triple $\mathcal{L} = (N, B, C)$ where
$N$ is a dictionary of \lstinline|IRNode| objects, $B$ is a binding
environment, and $C$ is a list of constraints.

\begin{definition}[IR Stack]\label{def:ir-stack}
The \emph{IR Stack} is an ordered sequence
$\IRSt = (\mathcal{L}_0, \mathcal{L}_1, \ldots, \mathcal{L}_n)$
of IR layers connected by \emph{lowering passes}
$\Lower_{k \to k+1} : \mathcal{L}_k \to \mathcal{L}_{k+1}$.
Each pass satisfies the \emph{ambiguity-preservation invariant}:
if $\mathcal{L}_k$ contains an ambiguity mark $\mu$,
then $\mathcal{L}_{k+1}$ also contains $\mu$.
\end{definition}

The six IR layer kinds and their roles are:

\begin{center}
\begin{tabular}{lll}
  \toprule
  Layer kind & Symbol & Description \\
  \midrule
  \lstinline|SURFACE|       & $\mathcal{L}_0$ & Raw parsed term structure \\
  \lstinline|SEMANTIC|      & $\mathcal{L}_1$ & Desugared, name-resolved \\
  \lstinline|LOGICAL|       & $\mathcal{L}_2$ & Obligations extracted; predicates inlined \\
  \lstinline|SOLVER_READY|  & $\mathcal{L}_3$ & Fully encoded; ready for Z3 \\
  \lstinline|CACHED|        & $\mathcal{L}_4$ & Previously computed results \\
  \lstinline|DELTA|         & $\mathcal{L}_5$ & Incremental changes vs.\ prior run \\
  \bottomrule
\end{tabular}
\end{center}

\subsection{Deduction rules on the IR Stack}

Rules in the \textsc{DRE} operate at the $\mathcal{L}_2$ (Logical)
layer, where obligations have been extracted and predicates inlined.
The result of a rule application is a new node added to $\mathcal{L}_2$
and propagated forward to $\mathcal{L}_3$ by the subsequent lowering
pass.

\begin{proposition}[Lowering Commutes with Deduction]\label{prop:lower-commute}
  Let $r$ be any primitive rule and let $\Lower_{2 \to 3}$ be the
  lowering pass from Logical to Solver-Ready.  Then:
  \[
    \Lower_{2 \to 3}(\Apply(r, \mathcal{L}_2)) =
    \Apply(\Lower(r), \Lower_{2 \to 3}(\mathcal{L}_2))
  \]
  where $\Lower(r)$ is the solver-ready encoding of rule $r$.
\end{proposition}

This commutativity property is critical: it allows the \textsc{DRE} to
perform deduction at the logical level and have the results
automatically propagated to the solver level, without re-running
deduction on the lower representation.

\subsection{Normal forms and layer merging}

The fully flattened stack is defined as
\[
  \bigoplus_{k=0}^{n} \mathcal{L}_k
\]
where $\oplus$ is the \emph{layer merge operator}: right-biased union on
bindings, and union on nodes and constraints.  The normal form of a
proof tree (§\ref{sec:chaining}) corresponds to the merged stack with
all redundant constraint duplicates removed.

\begin{lstlisting}[style=jugeo-python,
  caption={Applying a deduction rule at the Logical IR layer.}]
from jugeo.encodings.ir_stack.lowering import LoweringPass
from jugeo.encodings.ir_stack.models import IRLayerKind

# Locate the logical layer
logical_layer = ir_stack.get_layer(IRLayerKind.LOGICAL)

# Apply conjunction-introduction at the logical level
result = apply_deduction_rule(
    rule    = conj_intro,
    judgment = logical_layer.get_judgment(goal_id),
    context  = {"available_judgments": logical_layer.all_judgments()},
)

# The lowering pass propagates the result to SOLVER_READY
solver_layer = LoweringPass(IRLayerKind.LOGICAL, IRLayerKind.SOLVER_READY
                            ).run(logical_layer.with_result(result))
\end{lstlisting}

% =====================================================================
\section{Integration with the Descent Engine}\label{sec:descent}
% =====================================================================

\subsection{Deduction as local section construction}

In the sheaf-theoretic framework of Paper~3, a global proof corresponds
to a \emph{global section} of the evidence sheaf $\evid$ over the
semantic site $\Site$.  The Descent Engine constructs global sections by
\emph{gluing} compatible local sections across covering families.

The role of \textsc{DRE} is to supply these local sections.  Given a
covering family $\{U_i\}$ of a coordinate $c$, the rule engine
processes each restriction $\J|_{U_i}$ independently, building a local
proof tree $\PTree_i : \Prov{\Gamma_i}{\phi|_{U_i}}$.  These local
trees are then submitted to the Descent Engine for global gluing.

\begin{definition}[Local Section from Deduction]\label{def:local-section}
A \emph{local section} produced by \textsc{DRE} at an open $U_i$ is a
pair $(\PTree_i, \tau_i)$ where $\PTree_i$ is a valid proof tree for
$\phi|_{U_i}$ and $\tau_i \in \Talg$ is the trust annotation of the
proof.  Local sections are \emph{compatible} if, for every overlap
$U_{ij} = U_i \cap U_j$, the restrictions satisfy
\[
  \res_{ij}(\PTree_i) = \res_{ji}(\PTree_j)
\quad \text{and} \quad
  \tau_i|_{U_{ij}} = \tau_j|_{U_{ij}}.
\]
\end{definition}

\subsection{The gluing pipeline}

The full pipeline from local evidence to global certificate proceeds as:

\begin{enumerate}[leftmargin=2em]
  \item \textbf{Decompose.}  The coordinator decomposes the goal
    judgment $\J$ into local sub-judgments $\J|_{U_i}$ over a
    covering family $\{U_i\}$.
  \item \textbf{Deduce.}  \textsc{DRE} applies rules to produce a local
    proof tree $\PTree_i$ for each $\J|_{U_i}$.
  \item \textbf{Check compatibility.}  The Descent Engine checks that
    the local sections satisfy the overlap conditions of
    Definition~\ref{def:local-section}.
  \item \textbf{Glue.}  The Descent Engine constructs the global section
    $\PTree : \Prov{\Gamma}{\phi}$ by the universal property of the
    sheaf.
  \item \textbf{Annotate.}  The global trust tier is computed as
    $\tau = \bigoplus_i \tau_i$ where $\tjoin$ is the conservative
    join of the trust algebra.
\end{enumerate}

\subsection{Trust tier through descent}

The trust tier of the global certificate is bounded by the trust of the
weakest local proof.  This reflects the principle that a chain of
deductions is only as strong as its weakest link.  In practice, the
bottleneck is usually the tier of the SMT solver oracle ($\Tsolver$)
when local proofs involve non-trivial arithmetic.

\begin{proposition}[Trust Conservation under Gluing]\label{prop:trust-glue}
  Let $\{(\PTree_i, \tau_i)\}$ be a compatible family of local sections
  produced by \textsc{DRE}.  Then the glued global certificate has trust
  tier $\tau_{\mathit{global}} = \bigoplus_i \tau_i$ where $\tjoin$ is
  the conservative join (Definition~3.2 of Paper~4).
\end{proposition}

% =====================================================================
\section{The Deduction Soundness Theorem}\label{sec:soundness}
% =====================================================================

\subsection{Formal setting}

We define a classical propositional semantics and prove that every
proof tree produced by the rule engine encodes a valid deduction.

\begin{definition}[Valuation and Satisfaction]\label{def:valuation}
A \emph{valuation} $\Val$ is a function $\Val : \mathsf{Atom} \to
\{\mathtt{true}, \mathtt{false}\}$.  The \emph{semantic bracket}
$\sem{\phi}_\Val$ is defined inductively:
\begin{align*}
  \sem{a}_\Val &= \Val(a) \\
  \sem{\phi \wedge \psi}_\Val &= \sem{\phi}_\Val \wedge
                                  \sem{\psi}_\Val \\
  \sem{\phi \to \psi}_\Val &=
    \neg\sem{\phi}_\Val \vee \sem{\psi}_\Val \\
  \sem{\top}_\Val &= \mathtt{true}
\end{align*}
A valuation $\Val$ \emph{satisfies} $\Gamma$, written $\Val \Sat \Gamma$,
if $\sem{\phi}_\Val = \mathtt{true}$ for every $\phi \in \Gamma$.
\end{definition}

\begin{definition}[Semantic Validity]\label{def:validity}
A sequent $\Prov{\Gamma}{\phi}$ is \emph{semantically valid} if
for every valuation $\Val$:
\[
  \Val \Sat \Gamma \implies \sem{\phi}_\Val = \mathtt{true}.
\]
\end{definition}

\subsection{Rule-by-rule soundness}

We first verify that each primitive rule is individually sound.

\begin{lemma}[Soundness of $\WeakR$]\label{lem:weak-sound}
  If $\Prov{\Gamma}{\phi}$ is semantically valid and $A$ is any
  formula, then $\Prov{\Gamma, A}{\phi}$ is semantically valid.
\end{lemma}

\begin{proof}
  Let $\Val \Sat \Gamma, A$.  Then $\Val \Sat \Gamma$, so
  $\sem{\phi}_\Val = \mathtt{true}$ by hypothesis.
\end{proof}

\begin{lemma}[Soundness of $\MP$]\label{lem:mp-sound}
  If $\Prov{\Gamma}{\phi \to \psi}$ and $\Prov{\Gamma}{\phi}$ are
  semantically valid, then so is $\Prov{\Gamma}{\psi}$.
\end{lemma}

\begin{proof}
  Let $\Val \Sat \Gamma$.  By hypothesis,
  $\sem{\phi \to \psi}_\Val = \mathtt{true}$
  and $\sem{\phi}_\Val = \mathtt{true}$.
  From $\sem{\phi \to \psi}_\Val = \neg\sem{\phi}_\Val \vee
  \sem{\psi}_\Val = \mathtt{true}$ and $\sem{\phi}_\Val = \mathtt{true}$,
  we obtain $\sem{\psi}_\Val = \mathtt{true}$.
\end{proof}

\begin{lemma}[Soundness of $\ConjI$]\label{lem:conj-i-sound}
  If $\Prov{\Gamma}{\phi}$ and $\Prov{\Gamma}{\psi}$ are semantically
  valid, then so is $\Prov{\Gamma}{\phi \wedge \psi}$.
\end{lemma}

\begin{lemma}[Soundness of $\ConjE$]\label{lem:conj-e-sound}
  If $\Prov{\Gamma}{\phi \wedge \psi}$ is semantically valid, then so
  are $\Prov{\Gamma}{\phi}$ and $\Prov{\Gamma}{\psi}$.
\end{lemma}

\subsection{Main theorem}

\begin{theorem}[Deduction Soundness]\label{thm:soundness}
  Let $\PTree : \Prov{\Gamma}{\phi}$ be a valid proof tree constructed
  by the Deduction Rule Engine using only primitive rules from the
  catalog of §\ref{sec:rules}.  Then the sequent
  $\Prov{\Gamma}{\phi}$ is semantically valid: for every valuation
  $\Val$, $\Val \Sat \Gamma$ implies $\sem{\phi}_\Val = \mathtt{true}$.
\end{theorem}

\begin{proof}
  By structural induction on the proof tree $\PTree$.

  \emph{Base case.}  If $\PTree$ is a leaf labelled by $\phi \in
  \Gamma$, then $\Val \Sat \Gamma$ directly gives
  $\sem{\phi}_\Val = \mathtt{true}$.

  \emph{Inductive step.}  Suppose $\PTree$ has root labelled by rule
  $r$ with children $\PTree_1, \ldots, \PTree_n$.  By the induction
  hypothesis, each $\PTree_i : \Prov{\Gamma_i}{\phi_i}$ is
  semantically valid.  We case-split on $r$:

  \begin{itemize}
    \item $r = \WeakR$: the sole premise is $\Prov{\Gamma}{\phi}$
      (valid by IH); Lemma~\ref{lem:weak-sound} gives validity of
      the conclusion.
    \item $r = \MP$: premises $\Prov{\Gamma}{\phi \to \psi}$ and
      $\Prov{\Gamma}{\phi}$ are both valid by IH;
      Lemma~\ref{lem:mp-sound} gives validity of $\Prov{\Gamma}{\psi}$.
    \item $r = \ConjI$: Lemma~\ref{lem:conj-i-sound} applies.
    \item $r = \ConjE_1$ or $r = \ConjE_2$: Lemma~\ref{lem:conj-e-sound}
      applies.
    \item $r = \UnivI$: universal instantiation is classically valid:
      $\forall x.\,\phi(x)$ is true under $\Val$ iff $\phi(t)$ is true
      under every extension of $\Val$ that assigns a value to $t$.
    \item $r = \FrameR$: the frame predicate $R$ is independent of
      the proof's conclusion by the side condition; the result follows
      from the IH and the semantics of $*$.
    \item $r = \SubstR$: substitution is classically valid by the
      substitution lemma.
    \item $r = \CutR$: by Theorem~\ref{thm:cut-elim}, every use of
      cut can be eliminated; the cut-free proof is covered by the
      remaining cases.
    \item $r = \mathsf{exchange}$, $\mathsf{contraction}$:
      classical tautologies.
  \end{itemize}

  In every case the conclusion is semantically valid.
\end{proof}

\subsection{Corollaries}

\begin{corollary}[Consistency]\label{cor:consistency}
  The rule engine never produces a proof of $\bot$ (falsehood) from a
  consistent context.
\end{corollary}

\begin{corollary}[Trust Soundness]\label{cor:trust}
  Let $(\PTree, \tau)$ be a trust-annotated proof tree produced by
  \textsc{DRE}.  If every leaf carries a trust tier $\geq \Tsolver$,
  then the root carries a trust tier $\geq \Tsolver$.
\end{corollary}

\begin{proof}
  By the trust-impact rules of §\ref{sec:rules}, the trust tier of
  each rule application is the $\tjoin$-minimum of its premises.
  Since $\tjoin$ is the conservative join (Paper~4, Definition~3.1),
  the trust of the root cannot exceed the trust of any leaf.  The
  $\Tsolver$ tier is the minimum tier for premises of type
  \lstinline|SEMANTIC|; structural rules do not change trust.
\end{proof}

\begin{corollary}[Schema Soundness]\label{cor:schema-sound}
  Every proof tree produced by instantiating and discharging a
  theorem schema $\Schema$ is sound: the instantiated statement
  $\Inst(\Schema, \theta).\mathit{stmt}$ is semantically valid.
\end{corollary}

% =====================================================================
\section{Experiments}\label{sec:experiments}
% =====================================================================

We evaluated \textsc{DRE} on the three benchmark suites from the
JuGeo evaluation framework.

\subsection{Setup}

All experiments were run on the same hardware as earlier papers in
this series.  We use the standard $\numcases$ benchmark cases split
evenly across the specification, equivalence, and bug-detection suites.
For each case, the rule engine is invoked with the obligation set
$\oblig$ extracted from the judgment and the evidence bundle $\evid$
as the initial fact set.

\subsection{Results}

\begin{table}[h]
\centering
\begin{tabular}{lrrrr}
  \toprule
  Suite & Cases & Resolved & Avg.\ chain & Avg.\ time \\
  \midrule
  Specification     & \numspec  & — & — & — \\
  Equivalence       & \numeq    & — & — & — \\
  Bug detection     & \numbug   & — & — & — \\
  \midrule
  Total / average   & \numcases & \expAccuracy & — & \subSiteMeanTime \\
  \bottomrule
\end{tabular}
\caption{Deduction Rule Engine performance across benchmark suites.
  ``Resolved'' means the engine produced a complete proof tree without
  escalating to the Descent Engine.  ``Avg.\ chain'' is the mean
  number of rule applications in the shortest valid proof tree.}
\label{tab:results}
\end{table}

\paragraph{Rule distribution.}
Figure~\ref{fig:rules} shows the distribution of rules across all
resolved proof trees.

\begin{figure}[h]
\centering
\begin{tabular}{lr}
  \toprule
  Rule & \% of applications \\
  \midrule
  $\MP$        & — \\
  $\ConjI$     & — \\
  $\WeakR$     & — \\
  $\ConjE$     & — \\
  $\SubstR$    & — \\
  $\UnivI$     & — \\
  $\FrameR$    & — \\
  \bottomrule
\end{tabular}
\caption{Distribution of rule applications across all resolved proofs.}
\label{fig:rules}
\end{figure}

\paragraph{Schema acceleration.}
Without the theorem schema library, the average chain length was
longer chains and more cases required manual rule suggestion.
With schemas, the average chain length decreases and manual
intervention is substantially reduced.

\paragraph{Unresolved cases.}
The remaining unresolved cases (across all suites) are escalated to
the Descent Engine, which uses the global gluing procedure to complete
the proof.  Of these, most are resolved by the Descent Engine
within $\totaltime$ total time, yielding an end-to-end resolution
rate approaching \expAccuracy.

\paragraph{Trust tier distribution.}
Of resolved proofs, most achieve $\Tsolver$ tier, some
achieve $\Tcopilot$ tier (where at least one leaf was a Copilot-oracle
inference), and a few remain at $\Tunver$ (insufficient local
evidence for a higher tier).

% =====================================================================
\section{Related Work}\label{sec:related}
% =====================================================================

\paragraph{Sequent calculi and rule engines.}
The foundational treatment of sequent calculi is Gentzen's original
work~\cite{gentzen1935}.  Modern automated theorem provers such as
Leanprover~\cite{moura2021lean4} and Coq~\cite{coq2024} implement
rule engines with tactic languages.  \textsc{DRE} differs in being
designed for the \emph{incremental} setting, where evidence arrives
online and rules must fire in bounded time.

\paragraph{Proof-carrying code.}
Proof-carrying code~\cite{necula1997} requires that programs include
correctness certificates.  The Deduction Rule Engine generates these
certificates compositionally from local evidence, complementing the
proof-carrying Python approach of Paper~9 in this series.

\paragraph{Inference in program verification.}
Dependent type systems such as $F^\star$~\cite{swamy2016dependent} and
Dafny~\cite{leino2010dafny} perform inference via bidirectional type
checking, which is a special case of backward chaining.  \textsc{DRE}
generalises this to the full JuGeo rule catalog, including the frame
and substitution rules needed for separation-logic obligations.

\paragraph{Judgment Geometry series.}
This paper extends the judgment algebra of Paper~2~\cite{jugeo02} with
a formal deduction system, the descent framework of
Paper~3~\cite{jugeo03} with a local section generator, the trust algebra
of Paper~4~\cite{jugeo04} with a compositional trust-propagation account,
and the IR Stack of Paper~32 with deduction-level operations.

% =====================================================================
\section{Conclusion}\label{sec:conclusion}
% =====================================================================

We have presented the Deduction Rule Engine, the component of \JuGeo{}
responsible for converting local evidence into structured proof trees.
The engine's nine primitive rules cover all reasoning patterns needed
by the JuGeo verification pipeline: structural manipulation, logical
decomposition, universal instantiation, frame separation, and variable
substitution.

Our main result---the Deduction Soundness Theorem---guarantees that
every proof tree produced by the engine is semantically valid.  The
proof proceeds by structural induction over proof trees, with
Lemmas~\ref{lem:weak-sound}--\ref{lem:conj-e-sound} covering the
individual rules.  Corollaries establish consistency, trust soundness,
and schema soundness.

The theorem schema library reduces average proof chain length by
significantly and nearly eliminates the need for manual rule suggestion.
The IR Stack integration ensures that deduction results are
automatically propagated from the logical layer to the solver-ready
layer, and the Descent Engine integration enables local proofs to be
glued into global certificates.

Future work includes:
(i) extending the rule catalog to handle higher-order
goals arising from recursive Python programs,
(ii) connecting the proof-tree representation to the Lean 4
formalisation for certificate extraction,
(iii) adding a resource-sensitive frame rule that tracks linearity, and
(iv) integrating the schema library with the Copilot oracle for
automated schema discovery.

% =====================================================================
% References
% =====================================================================
\begin{thebibliography}{99}

\bibitem{gentzen1935}
G.~Gentzen.
\newblock Untersuchungen \"uber das logische Schlie{\ss}en.
\newblock {\em Mathematische Zeitschrift}, 39(1):176--210, 1935.

\bibitem{necula1997}
G.~C.~Necula.
\newblock Proof-carrying code.
\newblock {\em Proc.\ POPL 1997}, pp.~106--119.

\bibitem{leino2010dafny}
K.~R.~M.~Leino.
\newblock Dafny: An automatic program verifier for functional correctness.
\newblock {\em Proc.\ LPAR 2010}, LNCS 6355, pp.~348--370.

\bibitem{moura2021lean4}
L.~de Moura and S.~Ullrich.
\newblock The {Lean~4} theorem prover and programming language.
\newblock {\em Proc.\ CADE 2021}, LNAI 12699.

\bibitem{swamy2016dependent}
N.~Swamy et~al.
\newblock Dependent types and multi-monadic effects in {F$^{\star}$}.
\newblock {\em Proc.\ POPL 2016}.

\bibitem{coq2024}
The Coq Development Team.
\newblock {\em The Coq Proof Assistant Reference Manual}.
\newblock Version 8.19, 2024.

\bibitem{jugeo02}
H.~Young.
\newblock Judgment algebra: An 8-tuple framework.
\newblock {\em Judgment Geometry, Paper~2}, 2024.

\bibitem{jugeo03}
H.~Young.
\newblock Sheaf descent and obstruction classes.
\newblock {\em Judgment Geometry, Paper~3}, 2024.

\bibitem{jugeo04}
H.~Young.
\newblock The trust ordered algebra.
\newblock {\em Judgment Geometry, Paper~4}, 2024.

\bibitem{jugeo09}
H.~Young.
\newblock Proof-carrying Python.
\newblock {\em Judgment Geometry, Paper~9}, 2024.

\bibitem{jugeo32}
H.~Young.
\newblock IR representations and the IR stack.
\newblock {\em Judgment Geometry, Paper~32}, 2024.

\bibitem{jugeo33}
H.~Young.
\newblock Theorem schema encoding.
\newblock {\em Judgment Geometry, Paper~33}, 2024.

\end{thebibliography}

\end{document}
